\section{Jump-measure design}
Jump banks are first-class experiment objects. The core banks are a double-well shell, complete graphs over minima, sparse edge graphs over minima (the triple-well adjacent and overlong supports), displacements lifted from latent minima through a linear mixing map (10D transformed M\"uller--Brown), random matched-length controls, and block flips for ManyWell. The random matched-length control preserves the empirical jump-length distribution of a reference graph while randomizing directions, so it tests whether physically meaningful geometry matters beyond nonlocality alone.

The comparison spine is local Langevin, raw compound-Poisson (CP), and the L\'evy-score-corrected CP (LSC-CP). Raw CP uses the identical jump law as LSC-CP but omits the stationary drift; sharing one random-number stream, the two processes differ only by the correction. Raw CP is therefore the central diagnostic---it exposes the invariant-measure defect that the L\'evy score removes---rather than a fair efficiency baseline. A separate efficiency benchmark adds MALA, BAOAB, HMC, and parallel tempering, with ESS reported alongside target-fidelity error.
